% ---------------------------------------------------------------------
\section{$\order{\as}$ 系数函数的推导}\label{sec:disall-nlo}

本节在 $D=4-2\varepsilon$ 维的常规维数正规化中计算夸克靶与胶子靶的部分子结构函数至 $\order{\as}$（图见\cref{fig:dis-oneloop}），并由 $\MSbar$ 质量因子化得到系数函数：
软—共线双极点与末态共线极点在实、虚修正之和中相消，剩下的初态共线极点恰好由 $\MSbar$ 部分子 PDF 吸收——这正是\cref{thm:disall-factorization}的机制在单圈的实现。

\begin{proposition}[$O(\alpha_s)$ 的 $\overline{\rm MS}$ 系数函数 \tagR]\label{prop:disall-nlo}
（$R_{\mathrm{PT}}$：在重整化微扰论中对红外正规化的部分子态严格。）以下陈述在单圈（$\order{\as}$）。（本命题只涉及部分子态。）
在常规维数正规化、入射胶子按 $D-2$ 个极化平均的 $\MSbar$ 方案中，按
\begin{equation}\label{eq:disall-norm}
\frac{F_2}{x}=\sum_qe_q^2\Big\{q+\bar q+\frac{\as}{2\pi}\big[C_2^q\otimes(q+\bar q)+2C_2^g\otimes g\big]\Big\},\qquad
\frac{F_L}{x}=\sum_qe_q^2\frac{\as}{2\pi}\big[C_L^q\otimes(q+\bar q)+2C_L^g\otimes g\big]
\end{equation}
归一化系数函数（$\sum_q$ 只对味求和；胶子项的因子 $2$ 计入 $\gamma^*g\to q\bar q$ 中同时出现的夸克与反夸克；$(C\otimes f)(x)=\int_x^1\frac{\dd\xi}\xi C(x/\xi)f(\xi)$），在 $\mu=Q$ 处
\begin{equation}\label{eq:disall-C}
\begin{aligned}
C_2^q(z)&=C_F\Big[2\mathcal D_1-\frac32\mathcal D_0-(1+z)\ln(1-z)-\frac{1+z^2}{1-z}\ln z\\
&\hspace{4.2em}+3+2z-\Big(\frac92+\frac{\pi^2}3\Big)\delta(1-z)\Big],\\
C_2^g(z)&=T_F\Big[\big(z^2+(1-z)^2\big)\ln\frac{1-z}{z}-8z^2+8z-1\Big],\\
C_L^q(z)&=2C_Fz,\qquad C_L^g(z)=4T_Fz(1-z),
\end{aligned}   % VF-DIS-7
\end{equation}
$\mathcal D_k\equiv[\ln^k(1-z)/(1-z)]_+$。对一般的 $\mu=\mu_F=\mu_R$，$C_2^i(z,Q^2/\mu^2)=C_2^i(z)+\ln(Q^2/\mu^2)P_{qi}(z)$（$i=q,g$），$C_L$ 在此阶与 $\mu$ 无关。
等价地 $C_2^q=C_F\big[\frac{1+z^2}{1-z}\big(\ln\frac{1-z}z-\frac34\big)+\frac{9+5z}4\big]_+$（加号作用于整个方括号）。若入射胶子按 2 个极化平均，则 $C_2^g\to C_2^g+P_{qg}$，$C_L^g$ 不变。
\uses{def:disall-breit,prop:sudakov-oneloop,lem:prelim-plus-expansion,prop:pdf-partonic-dr,lem:appC-phase-space,prop:dis-c-oneloop}
\end{proposition}

\begin{proof}
部分子结构函数 $\hat F^j_a/(ze_q^2)$（$j=q,g$）用\cref{def:disall-breit}的投影提取；$g_0^2=4\pi\as\mu^{2\varepsilon}S_\varepsilon^{-1}$，$S_\varepsilon=(4\pi)^\varepsilon\ee^{-\varepsilon\gamma_E}$；初态平均为夸克自旋 $\frac12$、颜色 $\frac1{N_c}$，胶子极化 $\frac1{D-2}$、颜色 $\frac1{N_c^2-1}$。
(1) Born：$\hat W^{\mu\nu}_{(0)}=\frac1{4\pi}2\pi\delta((p+q)^2)\frac12e_q^2\Tr[\sla p\gamma^\mu(\sla p+\sla q)\gamma^\nu]$；由 $\gamma^\mu\sla a\gamma_\mu=-(D-2)\sla a$、$\sla p\sla p=0$ 与
$\delta((p+q)^2)=\delta(1-z)/Q^2$，对任意 $D$ 有 $\hat F^{q(0)}_2/z=e_q^2\delta(1-z)$，$\hat F^{q(0)}_L=0$。   % VF-DIS-1
(2) 虚修正：由一圈形状因子（\cref{prop:sudakov-oneloop}），$\hat F^{q,\mathrm V}_2/(ze_q^2)=2F^{(1)}\delta(1-z)=\frac{\as}{2\pi}C_F(\frac{\mu^2}{Q^2})^\varepsilon[-\frac2{\varepsilon^2}-\frac3\varepsilon-8+\zeta_2+\order{\varepsilon}]\delta(1-z)$，极点全部是红外的（$Z_1=Z_2$，外线自能无标度），对 $F_L$ 无贡献。   % VF-DIS-3
(3) 实辐射 $\gamma^*q\to qg$：以 $D$ 维两体相空间（\cref{lem:appC-phase-space}）积分投影后的平方矩阵元，$y$ 积分均为 Beta 函数，对 $\varepsilon$ 精确地得到（中间步骤见\cref{ex:disall-nlo}）
\begin{equation}\label{eq:disall-Rq}
\begin{gathered}
\frac{\hat F^{q,\mathrm R}_a}{ze_q^2}=\frac{\as}{2\pi}C_F\Big(\frac{\mu^2}{Q^2}\Big)^\varepsilon R^q_a,\\
R^q_2=\frac{\ee^{\varepsilon\gamma_E}\Gamma(-\varepsilon)}{2\Gamma(2-2\varepsilon)}\frac{z^\varepsilon\,[2(1+z^2)-\varepsilon(7-2z)+\varepsilon^2]}{(1-z)^{1+\varepsilon}},\qquad
R^q_L=\frac{\ee^{\varepsilon\gamma_E}\Gamma(2-\varepsilon)}{\Gamma(2-2\varepsilon)}\,2z\Big(\frac z{1-z}\Big)^\varepsilon .   % VF-DIS-4
\end{gathered}
\end{equation}
(4) 用分布展开（\cref{lem:prelim-plus-expansion}，$a=1$）$(1-z)^{-1-\varepsilon}=-\frac1\varepsilon\delta(1-z)+\sum_{k\ge0}\frac{(-\varepsilon)^k}{k!}\mathcal D_k$ 与乘积约定，记 $\delta\equiv\delta(1-z)$，
\begin{equation}\label{eq:disall-Rq-exp}
\begin{aligned}
R^q_2={}&\frac{2\delta}{\varepsilon^2}+\frac1\varepsilon\Big[-2\mathcal D_0+1+z+\frac32\delta\Big]+2\mathcal D_1-\frac32\mathcal D_0\\
&-(1+z)\ln(1-z)-\frac{1+z^2}{1-z}\ln z+3+2z+\Big(\frac72-3\zeta_2\Big)\delta+\order{\varepsilon},   % VF-DIS-4
\end{aligned}
\end{equation}
而 $R^q_L=2z+\order{\varepsilon}$ 有限。(5) 实、虚相加：$1/\varepsilon^2$ 的系数 $2-2=0$；$1/\varepsilon$ 的系数 $-2\mathcal D_0+1+z+(\frac32-3)\delta=-P_{qq}/C_F$（$P_{qq}/C_F=2\mathcal D_0-(1+z)+\frac32\delta$）；
$\delta$ 的有限系数 $\frac72-3\zeta_2-8+\zeta_2=-(\frac92+\frac{\pi^2}3)$。于是 $\hat F^q_2/(ze_q^2)=\delta+\frac{\as}{2\pi}(\frac{\mu^2}{Q^2})^\varepsilon[-\frac1\varepsilon P_{qq}+C_2^q+\order{\varepsilon}]$。   % VF-DIS-6
(6) 质量因子化：由\cref{prop:pdf-partonic-dr}，$f_{q/q}=\delta-\frac{\as}{2\pi}\frac1\varepsilon P_{qq}+\order{\as^2}$；因子化公式用于夸克靶给出
$\frac{\as}{2\pi}C_2^q(z,Q^2/\mu^2)=\frac{\as}{2\pi}[(\frac{\mu^2}{Q^2})^\varepsilon(-\frac{P_{qq}}\varepsilon+C_2^q)+\frac{P_{qq}}\varepsilon]+\order{\varepsilon}=\frac{\as}{2\pi}[C_2^q+\ln\frac{Q^2}{\mu^2}P_{qq}]$；
胶子靶的极点由 $f_{q/g}+f_{\bar q/g}=-\frac{\as}{2\pi}\frac2\varepsilon P_{qg}$ 吸收（\cref{ex:disall-gluon}）。(7) 两种写法的 $C_2^q$ 在 $z<1$ 处逐点相同，且在 $[0,1]$ 上积分都为零（Adler 求和规则），故为同一分布；
按 2 个极化平均相当于把 $R^g_2$ 乘以 $1-\varepsilon$，极点项 $-\frac2\varepsilon P_{qg}$ 产生有限移动 $+2P_{qg}$，即 $C_2^g\to C_2^g+P_{qg}$。   % VF-CONV-6
这些结果的 Mellin 矩就是\cref{ch:dis}的单圈 Wilson 系数（\cref{prop:dis-c-oneloop}）。
\end{proof}

\noindent
换到 $a_s=\as/4\pi$ 的归一化（$c^{(1)}_{2,\mathrm{ns}}=2C_2^q$，$c^{(1)}_{2,g}=4n_fC_2^g$，$c^{(1)}_{L,\mathrm{ns}}=2C_L^q$，$c^{(1)}_{L,g}=4n_fC_L^g$，$T_F=\frac12$），上述结果与\cite[式~(4.1)--(4.4)、(B.3)--(B.4)]{Vermaseren:2005qc}
及\cite[式~(3)]{Moch:2004xu}逐项一致，那里归功于\cite{BardeenEtAl1978}（其式~(3.9) 去掉 $\ln4\pi-\gamma_E$ 项后恰为 $2C^q_{2,N}$）；参见\cite{FurmanskiPetronzio1982,EllisStirlingWebber1996}。

\begin{example}[$O(\alpha_s)$ 系数函数的完整推导]\label{ex:disall-nlo}
我们写出\cref{prop:disall-nlo}证明第 (3)(4) 步的中间计算。实辐射振幅（\cref{fig:dis-oneloop}）为
$M^\mu=e_qgt^a\bar u(p')[\gamma^\rho\frac{\sla p+\sla q}{\hat s}\gamma^\mu+\gamma^\mu\frac{\sla p-\sla k}{\hat t}\gamma^\rho]u(p)\varepsilon^*_\rho(k)$，$\hat s=Q^2(1-z)/z$，$\hat t=(p-k)^2$，$\hat u=(p-p')^2$，$\hat s+\hat t+\hat u=-Q^2$。
只有一个外胶子时 $k^\rho M_\rho=0$，极化求和可取 $-g_{\rho\sigma}$（我们也用物理极化求和验证了结果不变）。$D$ 维 Dirac 代数给出
$(-g_{\mu\nu})\sum M^\mu M^{\nu*}|_{\mathrm{Dirac}}=8(1-\varepsilon)[(1-\varepsilon)(-\frac{\hat t}{\hat s}-\frac{\hat s}{\hat t})+\frac{2\hat uQ^2}{\hat s\hat t}+2\varepsilon]$，$(p_\mu p_\nu)\sum M^\mu M^{\nu*}|_{\mathrm{Dirac}}=-4(1-\varepsilon)\hat u$，
再乘以 $e_q^2g_0^2$ 与色—自旋平均因子 $\frac12\cdot\frac1{N_c}\Tr(t^at^a)=\frac{C_F}2$。在 $\gamma^*$—部分子质心系中以 $y=(1-\cos\theta)/2$ 参数化（$\theta$ 为出射胶子与 $\vec p$ 的夹角），
$\hat t=-Q^2y/z$，$\hat u=-Q^2(1-y)/z$，$\int\dd\Phi_2=\frac1{8\pi}(\frac{4\pi}{\hat s})^\varepsilon\frac1{\Gamma(1-\varepsilon)}\int_0^1\dd y\,[y(1-y)]^{-\varepsilon}$，总前因子为
\begin{equation*}
\frac1{4\pi}\cdot\frac{C_F}2g_0^2\cdot\frac1{8\pi}\Big(\frac{4\pi}{\hat s}\Big)^\varepsilon\frac1{\Gamma(1-\varepsilon)}=\frac{\as}{2\pi}C_F\cdot\frac18\Big(\frac{\mu^2}{Q^2}\Big)^\varepsilon\frac{z^\varepsilon(1-z)^{-\varepsilon}\ee^{\varepsilon\gamma_E}}{\Gamma(1-\varepsilon)}.
\end{equation*}
$F_2$ 投影后的被积函数是 $8[(1-\varepsilon)y^2-2y(z-\varepsilon(1-z))+1+z^2-\varepsilon(1-z)^2]/[y(1-z)]+16(3-2\varepsilon)z(1-y)$。
记 $I_a=\int_0^1\dd y\,y^{a-\varepsilon}(1-y)^{-\varepsilon}$，则 $I_0=\Gamma^2(1-\varepsilon)/\Gamma(2-2\varepsilon)$，$I_1=I_0/2$，$I_{-1}=-(1-2\varepsilon)I_0/\varepsilon$；$y\to0$ 端点（胶子与入射夸克共线）的 $1/\varepsilon$ 来自 $I_{-1}$。
逐项积分并乘以 $(1-z)/(8I_0)$ 得 $-(1+z^2)/\varepsilon+(7-2z)/2-\varepsilon/2=-\frac1{2\varepsilon}[2(1+z^2)-\varepsilon(7-2z)+\varepsilon^2]$，连同前因子与 $\Gamma(-\varepsilon)=-\Gamma(1-\varepsilon)/\varepsilon$ 即得\cref{eq:disall-Rq}。
$F_L$ 投影后的被积函数 $32(1-\varepsilon)z(1-y)$ 在 $y\to0$ 处不奇异：$F_L$ 的投影压低了共线区域，$R^q_L$ 有限。

展开时写 $R^q_2=A(\varepsilon)(1-z)^{-1-\varepsilon}G(z,\varepsilon)$，$A=\frac{\ee^{\varepsilon\gamma_E}\Gamma(-\varepsilon)}{2\Gamma(2-2\varepsilon)}$，$G=z^\varepsilon[2(1+z^2)-\varepsilon(7-2z)+\varepsilon^2]$，$G(1,\varepsilon)=4-5\varepsilon+\varepsilon^2$。
于是 $\delta$ 的系数为 $-AG(1,\varepsilon)/\varepsilon=\frac2{\varepsilon^2}+\frac3{2\varepsilon}+\frac72-3\zeta_2+\order{\varepsilon}$，$\mathcal D_0$ 的系数为 $AG(1,\varepsilon)=-\frac2\varepsilon-\frac32+\order{\varepsilon}$，$\mathcal D_1$ 的系数为 $-\varepsilon AG(1,\varepsilon)=2+\order{\varepsilon}$，
正规部分 $A[G(z,\varepsilon)-G(1,\varepsilon)](1-z)^{-1-\varepsilon}=\frac{1+z}\varepsilon-(1+z)\ln(1-z)-\frac{1+z^2}{1-z}\ln z+3+2z+\order{\varepsilon}$，合起来即\cref{eq:disall-Rq-exp}。
最后，对 $N=1,\dots,8$ 以 40 位精度数值计算\cref{eq:disall-C}的 Mellin 矩，与\cref{ch:dis}的闭式一致到 $10^{-39}$；特别地 $C^q_{2,1}=0$（Adler 求和规则），   % VF-DIS-8
$C^q_{1,1}=C^q_{2,1}-C^q_{L,1}=-C_F$（Bjorken $\nu$ 求和规则；$C_1^i\equiv C_2^i-C_L^i$ 来自 $2F_1=F_2/x-F_L/x$）。
\uses{prop:disall-nlo}
\end{example}

\begin{example}[光子—胶子融合]\label{ex:disall-gluon}
$\gamma^*(q)+g(p)\to q(k_1)+\bar q(k_2)$：$M^\mu=e_qgt^a\bar u(k_1)[\gamma^\mu\frac{\sla k_1-\sla q}{\hat u}\gamma^\rho+\gamma^\rho\frac{\sla k_1-\sla p}{\hat t}\gamma^\mu]v(k_2)\varepsilon_\rho(p)$，
$\hat t=(p-k_1)^2=-Q^2y/z$，$\hat u=(p-k_2)^2=-Q^2(1-y)/z$（$\theta$ 为出射夸克与 $\vec p$ 的夹角；矩阵元在 $\hat t\leftrightarrow\hat u$ 下对称，\cref{fig:dis-oneloop}中按内线动量对两个传播子的命名不影响结果）。
投影后的迹为 $(-g_{\mu\nu})\sum M^\mu M^{\nu*}|_{\mathrm{Dirac}}=8(1-\varepsilon)[(1-\varepsilon)(\frac{\hat u}{\hat t}+\frac{\hat t}{\hat u})-\frac{2\hat sQ^2}{\hat t\hat u}-2\varepsilon]$，$(p_\mu p_\nu)\sum M^\mu M^{\nu*}|_{\mathrm{Dirac}}=8\hat s$，
平均因子为 $\frac1{2(1-\varepsilon)}\frac1{N_c^2-1}\Tr(t^at^a)=\frac{T_F}{2(1-\varepsilon)}$。记 $\hat F^g_a/(ze_q^2)=\frac{\as}{2\pi}T_F(\frac{\mu^2}{Q^2})^\varepsilon R^g_a$（一味夸克，$q$ 与 $\bar q$ 都计入），
$F_2$ 投影的 $y$ 积分为 $-\frac{16}\varepsilon\frac{\Gamma^2(1-\varepsilon)}{\Gamma(2-2\varepsilon)}\frac{P-3\varepsilon+4\varepsilon^2-2\varepsilon^3}{1-\varepsilon}$，$P\equiv z^2+(1-z)^2$，于是
\begin{equation*}
R^g_2=\frac{2\ee^{\varepsilon\gamma_E}\Gamma(-\varepsilon)}{\Gamma(2-2\varepsilon)}\Big(\frac z{1-z}\Big)^\varepsilon\frac{P-3\varepsilon+4\varepsilon^2-2\varepsilon^3}{(1-\varepsilon)^2}
=-\frac2\varepsilon P+2\Big\{P\ln\frac{1-z}z-8z^2+8z-1\Big\}+\order{\varepsilon},   % VF-DIS-5
\end{equation*}
而 $R^g_L=8z(1-z)+\order{\varepsilon}$。极点出现两次，$y\to0$ 与 $y\to1$（夸克或反夸克与胶子共线）各贡献 $-P/\varepsilon$。
因此 $\hat F^g_2/(ze_q^2)=\frac{\as}{2\pi}(\frac{\mu^2}{Q^2})^\varepsilon[-\frac2\varepsilon P_{qg}+2C_2^g+\order{\varepsilon}]$，减去 $f_{q/g}+f_{\bar q/g}=-\frac{\as}{2\pi}\frac2\varepsilon P_{qg}$ 与 Born 系数之卷积即得 $2[C_2^g+\ln(Q^2/\mu^2)P_{qg}]$。
按 2 个极化平均时 $R^g_2$ 乘以 $1-\varepsilon$，$C_2^g$ 移动 $+P_{qg}$：胶子平均的约定是 $\MSbar$ 方案的一部分。
\uses{prop:disall-nlo}
\end{example}

\noindent
Callan--Gross 关系 $F_L=0$ 只是部分子模型（领头阶）的结果，而不是 QCD 的定理：由 $C_L^q=2C_Fz$、$C_L^g=4T_Fz(1-z)$ 与领头阶的 $F_2/x=\sum e_q^2(q+\bar q)$ 立即得到
$F_L=\frac{\as}{4\pi}x^2\int_x^1\frac{\dd\xi}{\xi^3}\big[\frac{16}3F_2(\xi)+8\sum_qe_q^2(1-\frac x\xi)\xi g(\xi)\big]$（$C_F=\frac43$，$T_F=\frac12$），即 Altarelli--Martinelli 关系\cite{AltarelliMartinelli1978}；
$C_L^q$ 的矩最早由\cite{Zee:1974du}得到。$C_L$ 在 $\order{\as}$ 与 $\mu$ 及方案都无关，是胶子分布的直接检验之一。
